
\documentclass[11pt,a4paper]{article}
\usepackage[T1]{fontenc}
\usepackage{lmodern}
\usepackage{amsmath,amssymb,amsfonts,mathtools}
\usepackage{graphicx,booktabs,longtable,array,tabularx,adjustbox}
\usepackage{geometry,microtype,fancyhdr,lastpage,etoolbox,enumitem}
\usepackage{newunicodechar,textcomp}
\usepackage{hyperref}
\hypersetup{hidelinks,pdfcreator={XeLaTeX / QM1 academic edition}}
\geometry{left=20mm,right=20mm,top=22mm,bottom=23mm}
\setcounter{secnumdepth}{2}
\setlength{\parindent}{0pt}
\setlength{\parskip}{4pt plus 1pt}
\setlength{\emergencystretch}{2em}
\setlength{\tabcolsep}{3pt}
\renewcommand{\arraystretch}{1.16}
\setlist{nosep,topsep=3pt,itemsep=2pt,leftmargin=7mm}
\AtBeginEnvironment{longtable}{\footnotesize\setlength{\tabcolsep}{2.2pt}}
\newcommand{\tightlist}{\setlength{\itemsep}{0pt}\setlength{\parskip}{0pt}}
\providecommand{\Cov}{\operatorname{Cov}}
\providecommand{\Var}{\operatorname{Var}}
\providecommand{\Corr}{\operatorname{Corr}}
\providecommand{\MAD}{\operatorname{MAD}}
\providecommand{\E}{\mathbb{E}}
\providecommand{\R}{\mathbb{R}}
\newunicodechar{—}{\textemdash}
\newunicodechar{–}{\textendash}
\newunicodechar{€}{\texteuro}
\newunicodechar{≤}{\ensuremath{\leq}}
\newunicodechar{≥}{\ensuremath{\geq}}
\newunicodechar{×}{\ensuremath{\times}}
\newunicodechar{÷}{\ensuremath{\div}}
\newunicodechar{Σ}{\ensuremath{\Sigma}}
\newunicodechar{σ}{\ensuremath{\sigma}}
\newunicodechar{μ}{\ensuremath{\mu}}
\newunicodechar{Δ}{\ensuremath{\Delta}}
\newunicodechar{∑}{\ensuremath{\sum}}
\newunicodechar{≈}{\ensuremath{\approx}}
\newunicodechar{√}{\ensuremath{\surd}}
\newunicodechar{≠}{\ensuremath{\ne}}
\newunicodechar{→}{\ensuremath{\to}}
\newunicodechar{∞}{\ensuremath{\infty}}
\newunicodechar{∈}{\ensuremath{\in}}
\newunicodechar{−}{\ensuremath{-}}
\newunicodechar{²}{\ensuremath{^{2}}}
\newunicodechar{³}{\ensuremath{^{3}}}
\newunicodechar{₀}{\ensuremath{_{0}}}
\newunicodechar{₁}{\ensuremath{_{1}}}
\newunicodechar{₂}{\ensuremath{_{2}}}
\newunicodechar{₃}{\ensuremath{_{3}}}
\newunicodechar{₄}{\ensuremath{_{4}}}
\newunicodechar{ₜ}{\ensuremath{_{t}}}
\newunicodechar{ₚ}{\ensuremath{_{p}}}
\newunicodechar{ₙ}{\ensuremath{_{n}}}
\newunicodechar{ₓ}{\ensuremath{_{x}}}
\newunicodechar{ᵢ}{\ensuremath{_{i}}}
\newunicodechar{ŷ}{\ensuremath{\hat y}}
\newunicodechar{Ŷ}{\ensuremath{\widehat Y}}
\newunicodechar{ₘ}{\ensuremath{_{m}}}
\newunicodechar{ₑ}{\ensuremath{_{e}}}
\newunicodechar{ₐ}{\ensuremath{_{a}}}
\newunicodechar{ₖ}{\ensuremath{_{k}}}
\newunicodechar{ₛ}{\ensuremath{_{s}}}
\newunicodechar{ₗ}{\ensuremath{_{l}}}
\newunicodechar{ₕ}{\ensuremath{_{h}}}
\newunicodechar{ₒ}{\ensuremath{_{o}}}
\newunicodechar{₊}{\ensuremath{_{+}}}
\newunicodechar{₋}{\ensuremath{_{-}}}
\newunicodechar{₌}{\ensuremath{_{=}}}
\newunicodechar{₍}{\ensuremath{_{(}}}
\newunicodechar{₎}{\ensuremath{_{)}}}
\newunicodechar{⁻}{\ensuremath{^{-}}}
\newunicodechar{⁺}{\ensuremath{^{+}}}
\newunicodechar{⁰}{\ensuremath{^{0}}}
\newunicodechar{¹}{\ensuremath{^{1}}}
\newunicodechar{⁴}{\ensuremath{^{4}}}
\newunicodechar{⁵}{\ensuremath{^{5}}}
\newunicodechar{⁶}{\ensuremath{^{6}}}
\newunicodechar{⁷}{\ensuremath{^{7}}}
\newunicodechar{⁸}{\ensuremath{^{8}}}
\newunicodechar{⁹}{\ensuremath{^{9}}}
\makeatletter
\renewcommand{\section}{\@startsection{section}{1}{0pt}{2.0ex plus .7ex minus .3ex}{1.0ex}{\large\bfseries}}
\renewcommand{\subsection}{\@startsection{subsection}{2}{0pt}{1.8ex plus .5ex minus .2ex}{.7ex}{\normalsize\bfseries}}
\renewcommand{\subsubsection}{\@startsection{subsubsection}{3}{0pt}{1.3ex plus .4ex minus .2ex}{.5ex}{\normalsize\bfseries}}
\makeatother
\pagestyle{fancy}
\fancyhf{}
\renewcommand{\headrulewidth}{0.35pt}
\renewcommand{\footrulewidth}{0pt}
\fancyhead[L]{\small\textsc{KEDGE Business School --- QM1}}
\fancyhead[R]{\small\textsc{2026--2027}}
\fancyfoot[L]{\small Badr ABOUFARASSE, PhD}
\fancyfoot[R]{\small\thepage\ /\ \pageref{LastPage}}
\setlength{\headheight}{14pt}
\setkeys{Gin}{keepaspectratio}
\begin{document}

\begin{center}
{\small\textsc{KEDGE Business School \quad EBP1 \quad Quantitative Methods 1}}\\[6pt]
{\Large\bfseries Fiche 5 — Séries temporelles}\\[6pt]
{\normalsize\bfseries Énoncés des exercices}\\[6pt]
{\small Badr ABOUFARASSE, PhD \quad --- \quad 2026--2027}
\end{center}
\vspace{2mm}\hrule\vspace{6mm}

\hypertarget{applications-fondamentales}{%
\section{Applications fondamentales}\label{applications-fondamentales}}

Maîtriser les définitions et les calculs essentiels du cours.

\emph{Application du cours --- 10 min}

\hypertarget{exercice-1-composantes-et-moduxe8les}{%
\subsection{Exercice 1 -- Composantes et modèles}\label{exercice-1-composantes-et-moduxe8les}}

\begin{enumerate}
\item
  Associer à {\(T_t\)}, {\(C_t\)}, {\(S_t\)} ou {\(R_t\)} : tendance de long terme, récession, pic récurrent du quatrième trimestre, choc exceptionnel.
\item
  Écrire le modèle additif complet et le modèle multiplicatif complet.
\item
  Expliquer dans quel cas une saisonnalité multiplicative est plus plausible qu'une saisonnalité additive.
\end{enumerate}

\emph{Application du cours --- 20 min}

\hypertarget{exercice-2-moyennes-mobiles-et-centrage}{%
\subsection{Exercice 2 -- Moyennes mobiles et centrage}\label{exercice-2-moyennes-mobiles-et-centrage}}

On observe deux années de demande trimestrielle : {\[80,\ 100,\ 120,\ 140,\ 88,\ 110,\ 132,\
154.\]}

\begin{enumerate}
\item
  Décrire le motif principal de la série.
\item
  Calculer les cinq moyennes mobiles d'ordre 4.
\item
  Les centrer afin d'obtenir les moyennes mobiles centrées aux périodes 3 à 6.
\item
  Expliquer pourquoi une moyenne mobile d'ordre pair doit être recentrée.
\end{enumerate}

\hypertarget{pruxe9paration-aux-uxe9valuations}{%
\section{Préparation aux évaluations}\label{pruxe9paration-aux-uxe9valuations}}

Résoudre des exercices de contrôle continu et de partiel en autonomie.

\emph{Préparation CC / partiel --- 25 min}

\hypertarget{exercice-3-duxe9composition-multiplicative-et-pruxe9vision}{%
\subsection{Exercice 3 -- Décomposition multiplicative et prévision}\label{exercice-3-duxe9composition-multiplicative-et-pruxe9vision}}

\begin{longtable}[]{@{}lllll@{}}
\toprule\noalign{}
Année & Q1 & Q2 & Q3 & Q4 \\
\midrule\noalign{}
\endhead
\bottomrule\noalign{}
\endlastfoot
2005 & 82 & 98 & 120 & 150 \\
2006 & 91 & 111 & 137 & 171 \\
2007 & 104 & 126 & 154 & 194 \\
\end{longtable}

On donne la tendance estimée : {\[\widehat
T_t=6.343t+86.939.\]} Les coefficients saisonniers normalisés sont : {\[S_{Q1}=0.789,\quad S_{Q2}=0.903,\quad
S_{Q3}=1.052,\quad S_{Q4}=1.255.\]}

\begin{enumerate}
\item
  Interpréter la pente de la tendance.
\item
  Vérifier approximativement la normalisation des coefficients saisonniers.
\item
  Calculer les quatre prévisions trimestrielles pour 2008.
\item
  Expliquer comment désaisonnaliser une observation {\(Y_t\)} dans le modèle multiplicatif.
\end{enumerate}

\emph{Préparation CC / partiel --- 65 min}

\hypertarget{exercice-5-chauxeene-compluxe8te-de-duxe9composition-saisonniuxe8re-type-partiel}{%
\subsection{Exercice 5 --- Chaîne complète de décomposition saisonnière --- type partiel}\label{exercice-5-chauxeene-compluxe8te-de-duxe9composition-saisonniuxe8re-type-partiel}}

Une entreprise suit ses ventes trimestrielles pendant trois ans. Utiliser le modèle multiplicatif Yₜ=Tₜ×Sₜ, sans cycle ni composante irrégulière explicites. Ne pas arrondir les résultats intermédiaires.

\begin{longtable}[]{@{}lllll@{}}
\toprule\noalign{}
Année / Year & Q1 & Q2 & Q3 & Q4 \\
\midrule\noalign{}
\endhead
\bottomrule\noalign{}
\endlastfoot
2019 & 40 & 52 & 48 & 68 \\
2020 & 44 & 58 & 53 & 76 \\
2021 & 50 & 65 & 60 & 84 \\
\end{longtable}

\begin{enumerate}
\tightlist
\item
  Tracer la série et calculer les neuf moyennes mobiles d\textquotesingle ordre 4.
\item
  Centrer les moyennes mobiles (t=3 à 10). Construire le tableau t, CMA4, t², t×CMA4.
\item
  Estimer la tendance Tₜ=at+b par moindres carrés sur les huit CMA4.
\item
  Calculer Yₜ/Tₜ pour les 12 trimestres, les quatre moyennes saisonnières et leurs coefficients rectifiés de somme 4.
\item
  Prévoir les quatre trimestres de 2022 (t=13 à 16), en valeurs désaisonnalisées puis saisonnalisées.
\end{enumerate}

\hypertarget{approfondissement-facultatif}{%
\section{Approfondissement facultatif}\label{approfondissement-facultatif}}

Aller plus loin dans l'interprétation et les applications professionnelles.

\emph{Approfondissement --- 15 min}

\hypertarget{exercice-4-uxe9valuation-des-pruxe9visions}{%
\subsection{Exercice 4 -- Évaluation des prévisions}\label{exercice-4-uxe9valuation-des-pruxe9visions}}

Pour cet exercice, on suppose que les valeurs effectivement observées en 2008 ont été : {\[136,\ 162,\ 188,\
245.\]} À partir des prévisions de l'exercice 3 :

\begin{enumerate}
\item
  calculer les erreurs {\(e_t=Y_t-\hat Y_t\)} ;
\item
  calculer le MAE ;
\item
  calculer le RMSE ;
\item
  expliquer pourquoi une période de validation distincte de la période d'estimation est utile ;
\item
  rappeler pourquoi une prévision naïve constitue un benchmark utile.
\end{enumerate}

\emph{Approfondissement --- 20 min}

\hypertarget{compluxe9ment-1-validation-des-pruxe9visions-et-rupture-structurelle}{%
\subsection{Complément 1 --- Validation des prévisions et rupture structurelle}\label{compluxe9ment-1-validation-des-pruxe9visions-et-rupture-structurelle}}

La demande mensuelle observée vaut 100, 104, 108, 150, 154 et 158. Un modèle ajusté sur les trois premiers mois prévoit 100, 104, 108, 112, 116 et 120.

\begin{enumerate}
\tightlist
\item
  Calculez les erreurs des mois 4 à 6.
\item
  Calculez la MAE et l'erreur moyenne sur la validation.
\item
  Diagnostiquez le modèle.
\item
  Proposez une réponse responsable.
\end{enumerate}

\textbf{Indication méthodologique :}

Utilisez erreur = observé − prévu. Des erreurs toujours positives signalent une sous-prévision systématique.

\end{document}
